\begin{tikzpicture}[
    >=Stealth,
    stand/.style={line width=1.4pt, draw=black!70},
    baseplate/.style={draw=black!80, fill=black!20, rounded corners=1pt},
    sensor/.style={circle, draw=black!85, fill=gray!25, line width=1pt, minimum size=0.55cm},
    source/.style={circle, draw=black!85, fill=orange!75, line width=1pt, minimum size=0.55cm},
    tabletop/.style={draw=black!70, line width=0.8pt, pattern=north east lines, pattern color=black!25},
    raydir/.style={->, line width=1.1pt, draw=UdeAGreen!85!black},
    rayrefl/.style={->, line width=1pt, dashed, draw=red!70!black}
]

% 1. Superficie reflectora horizontal
\draw[tabletop] (-5.2,-0.35) rectangle (5.2,0.0);
\node[below=3pt, font=\footnotesize] at (0,-0.35) {Superficie reflectora horizontal (mesa o suelo)};

% 2. Posiciones físicas (elevadas para dar espacio)
\coordinate (Mic) at (-3.4, 2.2);
\coordinate (Spk) at (3.4, 2.2);
\coordinate (Bounce) at (0.0, 0.0);

% Soporte receptor
\draw[baseplate] (-3.8,0.0) rectangle (-3.0,0.12);
\draw[stand] (-3.4,0.12) -- (Mic);
\node[sensor] (NodeMic) at (Mic) {};
\node[left=6pt, font=\small\bfseries] at (Mic) {Receptor ($R$)};

% Soporte emisor
\draw[baseplate] (3.0,0.0) rectangle (3.8,0.12);
\draw[stand] (3.4,0.12) -- (Spk);
\node[source] (NodeSpk) at (Spk) {};
\node[right=6pt, font=\small\bfseries] at (Spk) {Fuente ($S$)};

% Frentes de onda esquemáticos
\draw[line width=0.8pt, orange!85!black] ($(Spk)+(-0.4,0.25)$) arc (150:210:0.5);
\draw[line width=0.8pt, orange!85!black] ($(Spk)+(-0.7,0.40)$) arc (150:210:0.8);

% 3. Trayectorias acústicas
\draw[raydir] (NodeSpk.west) -- node[above=2pt, font=\footnotesize] {Onda directa: $r_{\mathrm{dir}} = r$} (NodeMic.east);

\draw[rayrefl] (NodeSpk.south) -- node[below right=1pt, font=\footnotesize, text=red!70!black] {$r_{\mathrm{refl}}/2$} (Bounce);
\draw[rayrefl] (Bounce) -- node[below left=1pt, font=\footnotesize, text=red!70!black] {$r_{\mathrm{refl}}/2$} (NodeMic.south);
\fill[red!70!black] (Bounce) circle (1.8pt);

% Altura h
\draw[dashed, draw=black!40, line width=0.6pt] (0.0,0.0) -- (0.0,2.2);
\draw[<->, draw=black!70, line width=0.6pt] (-0.35,0.05) -- node[left=1pt, font=\small] {$h$} (-0.35,2.15);

% 4. Recuadro de fórmulas ubicado holgadamente ARRIBA del rayo directo
\node[
    draw=black!70,
    line width=0.7pt,
    fill=white,
    rounded corners=3pt,
    inner sep=6pt,
    anchor=south
] at (0.0, 3.1) {
    \footnotesize
    \begin{tabular}{@{} l l @{}}
        \multicolumn{2}{@{} l @{}}{\textbf{Geometría de reflexión y selección temporal:}} \\[3pt]
        Trayectoria directa: & $r_{\mathrm{dir}} = r$ \\
        Trayectoria reflejada: & $r_{\mathrm{refl}} = \sqrt{r^2+4h^2}$ \\
        Diferencial de camino: & $\Delta r_{\mathrm{refl}} = \sqrt{r^2+4h^2} - r$ \\[4pt]
        \multicolumn{2}{@{} l @{}}{\textcolor{UdeAGreen!85!black}{\textbf{Condición de selección temporal:}} \quad $c(T)T_{\mathrm{gate}} < \Delta r_{\mathrm{refl}}$}
    \end{tabular}
};

\end{tikzpicture}